\documentclass[aspectratio=169]{beamer}
\usepackage[utf8]{inputenc}\usepackage{amsmath}\usepackage{tcolorbox}
\definecolor{UNIBBlue}{RGB}{0, 51, 102}\setbeamercolor{frametitle}{bg=UNIBBlue,fg=white}
\title{Optimisasi untuk Teknik Elektro}\subtitle{Minggu 13: Economic Dispatch dengan PSO}
\date{Semester Ganjil 2026}\begin{document}\begin{frame}\titlepage\end{frame}
\begin{frame}{Economic Dispatch (ED)}
\textbf{Tujuan:} Menentukan daya output (MW) tiap generator agar biaya bahan bakar minimum dan beban sistem terpenuhi.
\textbf{Mengapa PSO?} Kurva biaya generator riil memiliki \textit{valve-point effects} (efek non-linier/non-convex ripple) yang membuat metode turunan klasik gagal menemukan optimum global. PSO sangat andal untuk mengatasi masalah ini.
\end{frame}\end{document}